\documentclass[conference]{IEEEtran}
\usepackage{gvv-book}
\usepackage{gvv}
\usepackage{algorithm}
\usepackage{algpseudocode}
\usepackage{listings}
\usepackage{xcolor}
\usepackage{amsmath}
\usepackage{graphicx}
\usepackage{float}
\usepackage{booktabs}
\usepackage{url}
\usepackage{circuitikz}
\usepackage{hyperref}
\usepackage{tabularx}

\lstset{
  basicstyle=\ttfamily\small,
  keywordstyle=\color{blue}\bfseries,
  commentstyle=\color{green!40!black},
  stringstyle=\color{red},
  xleftmargin=1.5em,
  frame=single,
  breaklines=true,
  breakatwhitespace=true,
  tabsize=2,
  showstringspaces=false,
  captionpos=b,
  keepspaces=true
}
\def\UrlBreaks{\do\/\do-\do\_\do\.}

\title{\textbf{Migrating Scientific Calculator on ATtiny-85 from ATmega-328p}}

\author{
\IEEEauthorblockN{EE25BTECH11030-Josyula Avaneesh \and EE25BTECH11054-Harsha Soma Vardhan}
}

\begin{document}

\maketitle

\textbf{Introduction:}\\
This report details the architectural decisions, hardware selection, and algorithmic optimizations required to implement a fully functional order-of-operations (Shunting Yard) scientific calculator on an ATtiny85 microcontroller. Constrained by 6,012 bytes of usable Flash memory and lacking native floating-point hardware, the project required transitioning to a custom 32-bit fixed-point mathematical engine, eliminating standard libraries, and aggressively optimizing UI elements to prevent memory overflow.

\section{Hardware Architecture and Cost Analysis}

\subsection{Bill of Materials}
The hardware was selected to balance cost, form factor, and input/output capabilities. The primary challenge was interfacing a 25-button matrix and a 16x2 LCD using the highly limited GPIO pins of the ATtiny85. The components were sourced using standard Indian electronics market pricing.

\begin{table}[htbp]
\centering
\caption{Bill of Materials and Project Cost (in INR)}
\resizebox{\columnwidth}{!}{%
\begin{tabular}{@{}llcc@{}}
\toprule
\textbf{Component} & \textbf{Purpose} & \textbf{Qty} & \textbf{Est. Cost} \\ \midrule
Digispark (ATtiny85) & Core MCU ($16\text{MHz}$, $8\text{KB}$ Flash) & 1 & 230Rs \\
16x2 LCD (with I2C) & Visual Display (reduces pin usage to 2) & 1 &  200Rs \\
74HC595,74HC165 Shift Registers & Matrix column/row scanning expansion & 1 each & 20Rs \\
Tactile Push Buttons & 5x5 Keypad Matrix & 25 & 50Rs \\
$220\Omega$, $10k\Omega$ Resistors & Current limiting / Short-circuit protection & 1, 5 & 10Rs \\
$100\text{nF}$ Capacitors & IC Decoupling \& Hardware Debouncing & 3 &10Rs \\
(Breadboard, wires) & Physical prototyping & 1 &  90Rs \\ \midrule
\textbf{Total} & & & \textbf{610Rs} \\ \bottomrule
\end{tabular}%
}
\end{table}

\subsection{Cost Comparison: ATtiny85 vs. ATmega328P}
An obvious architectural alternative to the ATtiny85 would have been an ATmega328P. The ATmega328P provides $32\text{KB}$ of Flash memory and sufficient native GPIO pins to drive the matrix directly without shift registers.

However, opting for an Atmega328p would increase the core microcontroller cost from 250Rs to approximately 400Rs - 500Rs. While the ATmega328P easily accommodates bloated standard C \texttt{<math.h>} floating-point libraries, it removes the strict memory constraints that forced the creative algorithmic problem-solving at the heart of this project, allowing it to also include functions like exponential and logarithmic, along with being able to use fundamental algorithms like RK4. The ATtiny85 was deliberately selected to keep the budget highly optimized while functioning as close to an actual scientific calculator as possible.

Upon calculating the total cost while using an Atmega328p instead of an ATtiny85, we find that the total is \textbf{800Rs}. This increase in price is due to the Atmega328p not having a built-in oscillator and being unable to regulate power.

\section{Circuit Operation and Component Justification}

\subsection{System Operation Overview}
The ATtiny85 has only 6 standard GPIO pins. One is tied to the Reset function, leaving 5 usable pins. To drive both an LCD and a 25-button matrix, extreme multiplexing is required:
\begin{itemize}
    \item \textbf{Pins PB0 \& PB2:} Dedicated to the I2C protocol (SDA/SCL) to communicate with the 16x2 LCD module.
    \item \textbf{Pins PB1, PB3, PB4:} Dedicated to a serial Shift Register protocol (Clock, Latch, Data) to sequentially scan the 5x5 keypad matrix. The Data pin (PB4) acts dynamically as both an output to push row states and an input to read column states.
\end{itemize}

\subsection{Justification of Passive Components}
\begin{itemize}
    \item \textbf{$10k\Omega$ Resistors:} Placed on the column lines of the button matrix. When pressing multiple buttons simultaneously in a standard matrix, a direct short circuit can occur between a HIGH output pin and a LOW output pin, potentially frying the microcontroller or the shift register. The $10k\Omega$ resistors safely limit this current to roughly $0.5\text{mA}$ at $5\text{V}$, protecting the ICs.
    \item \textbf{$100\text{nF}$ Ceramic Capacitors:} Used for logic-level decoupling. High-speed switching inside the ATtiny85 and the 74HC595 shift register causes micro-voltage dips on the power rail. Placing a $100\text{nF}$ capacitor physically close to the $V_{CC}$ and $GND$ pins of these ICs acts as a local energy reserve, smoothing out the voltage and preventing sudden microcontroller resets.
    \item \textbf{$220\Omega$ Shared Bus Hack:} To drive both the 74HC595 and 74HC165 shift registers using a single data pin (\texttt{PB4}), the chips must share a data bus, which introduces a critical hardware conflict known as bus contention,\textit{a conflict occurring when two or more devices or signals attempt to drive the same bus or node simultaneously, causing unpredictable outputs, increased power consumption, or potential hardware damage}. 
    
    Because the 74HC165's output pin (\texttt{Q7}) is always actively driving the line HIGH or LOW, a catastrophic dead short would occur if the ATtiny85 simultaneously output an opposing logic state during a write cycle to the 74HC595, easily destroying both ICs. Placing a $220\Omega$ resistor inline between \texttt{PB4} and \texttt{Q7} perfectly solves this by acting as a collision buffer. When a logic mismatch occurs (creating a $5\text{V}$ difference across the resistor), Ohm's Law ($I = 5\text{V} / 220\Omega$) dictates that the short-circuit current is safely throttled to approximately $22.7\text{mA}$. This limit sits comfortably below the $40\text{mA}$ absolute maximum rating of the ATtiny85 and the $\pm 25\text{mA}$ limit of the 74HC series, allowing the microcontroller to safely overpower the 165's signal to write data to the 595 while reading input flawlessly when the pin switches to a high-impedance state.
\end{itemize}

\section{Hardware Wiring Connections}

\begin{table*}[htbp]
\centering
\caption{Hardware Wiring Connections}
\renewcommand{\arraystretch}{1.2}
\begin{tabular}{|l|l|l|l|l|}
\hline
\textbf{Source Component} & \textbf{Source Pin} & \textbf{Target Component} & \textbf{Target Pin} & \textbf{Notes} \\
\hline
\hline
\multicolumn{5}{|c|}{\textbf{Power \& Decoupling}} \\
\hline
Power Supply & 5V & Global VCC Rail & All 5V Pins & Main logic power \\
\hline
Power Supply & GND & Global GND Rail & All GND Pins & Common ground \\
\hline
100nF Capacitor & Leg 1 & 74HC595 \& 165 & VCC (Pin 16) & Place near ICs \\
\hline
100nF Capacitor & Leg 2 & 74HC595 \& 165 & GND (Pin 8) & Decoupling \\
\hline
\hline
\multicolumn{5}{|c|}{\textbf{I2C LCD Display}} \\
\hline
Digispark & PB0 (SDA) & I2C LCD Module & SDA & I2C Data \\
\hline
Digispark & PB2 (SCL) & I2C LCD Module & SCL & I2C Clock \\
\hline
I2C LCD Module & VCC & Power Rail & 5V & LCD Power \\
\hline
I2C LCD Module & GND & Power Rail & GND & LCD Ground \\
\hline
\hline
\multicolumn{5}{|c|}{\textbf{Shift Register Control Bus}} \\
\hline
Digispark & PB1 (CLK) & 74HC595 & SHCP (Pin 11) & Shared Clock \\
\hline
Digispark & PB1 (CLK) & 74HC165 & CP (Pin 2) & Shared Clock \\
\hline
Digispark & PB3 (LATCH) & 74HC595 & STCP (Pin 12) & Shared Latch \\
\hline
Digispark & PB3 (LATCH) & 74HC165 & PL (Pin 1) & Shared Latch \\
\hline
\hline
\multicolumn{5}{|c|}{\textbf{Data Bus \& 220 Ohm Collision Protection}} \\
\hline
Digispark & PB4 (DATA) & 74HC595 & DS (Pin 14) & Direct serial write \\
\hline
Digispark & PB4 (DATA) & 220 Ohm Resistor & Leg 1 & Short-circuit buffer \\
\hline
220 Ohm Resistor & Leg 2 & 74HC165 & Q7/QH (Pin 9) & Protects PB4 \\
\hline
\hline
\multicolumn{5}{|c|}{\textbf{74HC595 (Row Driver - Output)}} \\
\hline
74HC595 & Q0 - Q4 & Keypad Matrix & Rows 1 - 5 & Drives rows LOW \\
\hline
74HC595 & OE (Pin 13) & Power Rail & GND & Output Enable \\
\hline
74HC595 & MR (Pin 10) & Power Rail & 5V & Master Reclear \\
\hline
\hline
\multicolumn{5}{|c|}{\textbf{74HC165 (Column Reader - Input)}} \\
\hline
74HC165 & D0 - D4 & Keypad Matrix & Cols 1 - 5 & Reads columns \\
\hline
74HC165 & CE (Pin 15) & Power Rail & GND & Clock Enable \\
\hline
74HC165 & SER (Pin 10) & Power Rail & GND & Serial Input (Tie low)\\
\hline
\hline
\multicolumn{5}{|c|}{\textbf{Matrix Pull-ups (10k Ohm)}} \\
\hline
Keypad Matrix & Cols 1 - 5 & 10k Ohm Resistors & Leg 1 & One per column \\
\hline
10k Ohm Resistors & Leg 2 & Power Rail & 5V & Pulls cols HIGH \\
\hline
\end{tabular}
\end{table*}

\section{Software Hardships and Explicit Compromises}

To fit a Shunting Yard mathematical evaluator into the remaining $\sim 6,000$ bytes of memory, severe software compromises were deliberately made:

\begin{enumerate}
    \item \textbf{Elimination of Floating-Point Hardware:} The standard \texttt{<math.h>} library ballooned the code to $7,610$ bytes. It was entirely replaced by a custom 32-bit fixed-point engine scaled by $1,000$.
    \item \textbf{Removal of Logarithms \& Floating Exponents:} Calculating natural logs ($\ln$) or exponents ($e^x$) using Taylor Series loops causes 32-bit integer overflow (exceeding $2.1$ billion) when using a $1000$-scale fixed-point system. These functions were abandoned in favor of integer-only powers.
    \item \textbf{UI Translation Layer Stripped:} To save $\sim 150$ bytes of memory for the Modulus and Absolute Value functions, the translation code that expanded inputs (e.g., converting 's' to "sin(") was deleted. The LCD now displays raw memory buffer inputs (e.g., \texttt{s30=}).
    \item \textbf{Forced Degree Mode:} Fixed-point math loses extreme precision when handling irrational numbers like radians ($\pi$). Trigonometry was hardcoded to Degrees to maintain perfect whole-number inputs for the lookup tables.
\end{enumerate}

\section{Code Efficiency and Flash Memory Map}
\begin{table}[htbp]
\centering
\caption{Estimated Flash Memory Allocation}
\resizebox{\columnwidth}{!}{%
\begin{tabular}{@{}lc@{}}
\toprule
\textbf{Software Module} & \textbf{Approx. Flash Footprint (Bytes)} \\ \midrule
Shunting Yard Evaluator (Order of Operations) & $\sim 1,800$ \\
Main Control Loop \& Matrix Scanner & $\sim 1,500$ \\
I2C Protocol \& LCD Drivers & $\sim 600$ \\
Fixed-Point Math Engine (Integer scaling, bounds) & $\sim 450$ \\
Trigonometric \texttt{PROGMEM} Lookup Tables & 224 \\
Compiler Overhead \& Vector Tables & $\sim 1,300$ \\ \midrule
\textbf{Total Utilized Memory} & \textbf{ 5,874 / 6,012 (97.7\%)} \\ \bottomrule
\end{tabular}%
}
\end{table}

\section{Mathematical Error Analysis: Fixed-Point vs. Continuous Algorithms}

Because the calculator utilizes a 1000-scale fixed-point engine and lookup tables, specific mathematical tolerances and truncation errors exist. To properly contextualize these errors, this system's accuracy is benchmarked against the standard floating-point algorithms typically employed on larger microcontrollers (such as the Arduino Uno/ATmega328P), specifically the 4th-order Runge-Kutta (RK4) method and the Quake III inverse square root algorithm.

\subsection{System-Specific Error Tolerances}
\begin{itemize}
    \item \textbf{General Truncation Error:} Division operations are integer-based. Any decimal precision beyond the third decimal place ($0.001$) is truncated, not rounded. Maximum general precision error is $-0.0009$.
    \item \textbf{Forward Trigonometry ($\sin$, $\cos$, $\tan$):} The system uses a 91-point integer lookup table mapped exactly to degrees ($0^\circ$ to $90^\circ$). Inputs are cast to integers before lookup. Therefore, $\sin(45.6^\circ)$ is evaluated as $\sin(45^\circ)$. The maximum angular resolution error is $0.99^\circ$.
    \item \textbf{Inverse Trigonometry ($\arcsin$, $\arccos$, $\arctan$):} Implemented via a 21-point lookup table with linear interpolation. Because the curve of $\arcsin(x)$ steepens dramatically near $x = 1$, the linear interpolation introduces an error margin of approximately $\pm 0.5^\circ$ for values between bounds.
    \item \textbf{Square Root:} Utilizes an integer-based Newton-Raphson approximation algorithm. Convergence is strictly integer-bounded, yielding a maximum floor truncation error of $-0.001$.
\end{itemize}

\subsection{Comparative Analysis: Table Lookups vs. ODE Solvers (RK4)}
Standard academic implementations of scientific calculators on microcontrollers calculate trigonometric and exponential functions by reformulating them as Ordinary Differential Equations (ODEs).

\begin{itemize}
    \item \textbf{Accuracy vs. Processing Cost:} The RK4 method achieves extreme floating-point accuracy, with its local truncation error bounded by $O(h^5)$, where $h$ is the step size. However, calculating four distinct variables ($k_1, k_2, k_3, k_4$) per step requires intense floating-point multiplication. Our ATtiny85 system abandons this continuous accuracy in favor of $O(1)$ constant-time memory lookups, trading $O(h^5)$ precision for a hardcoded worst-case resolution error of $0.99^\circ$.
    \item \textbf{Root Approximations:}Standard implementations on larger chips like the ATmega328P use the famous ``Quake III'' algorithm to calculate roots. This algorithm uses a specific bit-level ``magic constant'' (\texttt{0x5f3759df}) to rapidly guess the answer, but it fundamentally requires the system to support complex floating-point decimal numbers. Because our ATtiny85 completely removes the heavy floating-point library to save over $1.5\text{KB}$ of Flash memory, we cannot use this trick. Instead, our system uses a simpler, whole-number (integer) guessing method. While this integer method introduces a tiny precision error (truncating by up to $-0.001$), it successfully calculates roots without the massive memory overhead required by standard decimal algorithms.
\end{itemize}

Ultimately, the algorithms utilized in systems like the ATmega328p provide superior mathematical continuity and precision. However, they demand memory capacities far exceeding the ATtiny85's $6\text{KB}$ limit. The discrete fixed-point tables and bounded integer Newton-Raphson methods deployed in this project represent the necessary compromises required to force advanced scientific evaluation onto heavily constrained 8-bit architecture.

\section{Development and Debugging Code}
During the development phase, isolated code snippets were used to verify hardware functionality before integrating the complex math engine.

\subsection{LCD I2C "Blink" Debug Code}
This code was used to verify I2C communication and LCD addressing ($0\text{x}27$) without the overhead of matrix scanning.\\

The code is present in the $codes$ directory under the name $blink.c$

\subsection{Matrix Keypad Debug Code}
This snippet was utilized to confirm that the shift register (74HC595) was correctly pulsing the rows and that the ATtiny85 was successfully reading the columns, outputting the raw grid coordinate to the LCD.\\

The code is present in the $codes$ directory under the name $button\_matrix.c$

\section{Final Optimized Firmware}
The final production code incorporates the fixed-point math engine,the 91-point trigonometric lookup table, the 21-point inverse trigonometric lookup table, the complete Shunting Yard string parsing evaluator, and the dual-mode mapping system containing Modulo and Absolute Value functions.\\

The code is present in the $codes$ directory under the name $main.c$

\end{document}